\hypertarget{appendix-fe6b457d68d8}{}\label{referencestart:appendix-fe6b457d68d8}
\section{Logical versus bitwise operators}\label{appendix-fe6b457d68d8-logical-versus-bitwise-operators}

\textbf{Curriculum depth:} Introduced · \textbf{First seen:} C001 · \textbf{Latest development:} C001

\begin{allmightroledefinition}

\subsection{Definition and mechanism}\label{appendix-fe6b457d68d8-definition-and-mechanism}

Logical operators reduce operands to conditions; bitwise operators transform corresponding representation lanes.

\end{allmightroledefinition}

Built-in \texttt{\&\&}, \texttt{\textbar{}\textbar{}}, and \texttt{!} convert operands to truth conditions. \texttt{\&}, \texttt{\textbar{}}, \texttt{\^{}}, and \texttt{\textasciitilde{}} operate on integer representation bits. For 10 and 12, \texttt{10\ \&\ 12} is 8 whereas \texttt{10\ \&\&\ 12} is true. Overloaded operators may have additional semantics.

\subsection{Prerequisites}\label{appendix-fe6b457d68d8-prerequisites}

\begin{itemize}
\tightlist
\item
  \hyperref[appendix-24475e1596a8]{Boolean algebra} --- The logical domain is false/true.
\item
  \hyperref[appendix-68b52453737c]{Bitwise AND} --- Per lane: 0\&0=0, 0\&1=0, 1\&0=0, 1\&1=1.
\item
  \hyperref[appendix-a427ce8ce8a6]{Bitwise OR} --- Per lane: only 0\textbar0 is 0; the other three cases are 1.
\end{itemize}

\subsection{Used by}\label{appendix-fe6b457d68d8-used-by}

\begin{itemize}
\tightlist
\item
  \hyperref[appendix-2e245a6fa8d3]{Short-circuit evaluation} --- uses logical versus bitwise operators as part of its representation or reasoning.
\end{itemize}

\subsection{Recall and next depth}\label{appendix-fe6b457d68d8-recall-and-next-depth}

Explain logical versus bitwise operators without the label, work the example above, then state the assumption whose removal would change the result.

\textbf{Pending:} overloads, precedence, vector predicates, and branch lowering.

\subsection{Source and navigation}\label{appendix-fe6b457d68d8-source-and-navigation}

\hyperref[01-foundations-i-ch001-from-a-signal-to-an-interpretation]{C001 teaching chapter} · \hyperref[appendix-fd090d982a15]{Part definitions and mechanisms} · \href{../07-Vision-XR-and-Playable-Worlds/main.pdf\#nameddest=appendix-4ad47bdd269a}{Dependency Map (Vision, XR, and Playable Worlds, Reference notes: Dependency Map)}
